\documentclass[journal]{IEEEtran}

\usepackage[T1]{fontenc}
\usepackage{amsmath,amssymb,amsfonts}
\usepackage{booktabs}
\usepackage{tikz}
\usetikzlibrary{positioning,arrows.meta,fit,calc,shapes.geometric,decorations.pathreplacing}
\usepackage{pgfplots}
\pgfplotsset{compat=1.18}
\usepackage{multirow}
\usepackage{xcolor}
\usepackage{bm}
\usepackage{algorithm}
\usepackage{algorithmic}
\usepackage{graphicx}
\usepackage{url}
\usepackage{cite}
\usepackage{balance}

\newtheorem{proposition}{Proposition}

\definecolor{ssmcolor}{RGB}{66,133,244}
\definecolor{cnncolor}{RGB}{234,67,53}
\definecolor{sharedcolor}{RGB}{52,168,83}
\definecolor{newcolor}{RGB}{251,188,4}

\title{Bridging the Noise Robustness Gap Between SSMs and CNNs: SNR-Adaptive Architectures for Ultra-Compact Keyword Spotting}

\author{Jin~Ho~Choi%
\thanks{J.~H.~Choi is with SmartEar Company, Seoul, South Korea (e-mail: jinhochoi@smartear.co.kr).}%
\thanks{This paper extends the authors' preliminary work presented at Interspeech 2026~\cite{choi2026nanomamba_interspeech}. The extensions include: (1)~CNN structural advantage decomposition (Section~III-D), (2)~NC-SSM with per-sub-band selectivity, stationarity/spectral-flatness conditioning, and Learned Spectral Gate (Section~V), (3)~SAGN with SNR-Conditioned Channel Scale (Section~VI), (4)~extended evaluation with 6+ models across 36~conditions (Section~VIII), and (5)~cross-architecture analysis (Section~IX). Sections~III-D, V, and VI are entirely new; Sections~III-A--C and IV are substantially extended.}%
\thanks{Manuscript submitted March 2026.}}

\begin{document}
\maketitle

% ============================================================
% ABSTRACT
% ============================================================
\begin{abstract}
Selective state space models (SSMs) such as Mamba achieve strong accuracy for keyword spotting (KWS) on edge devices through input-dependent dynamics, but degrade more severely than CNNs under noise---a gap that persists even after noise-augmented training.
We identify \emph{multiplicative noise propagation} as the root cause: SSM input-dependent parameters ($\Delta$, $\mathbf{B}$, $\mathbf{C}$ derived from input) create $O(\sigma_n^6)$ output variance, versus $O(\sigma_n^2)$ for CNN fixed kernels---a 1{,}000$\times$ gap at $-15$\,dB.
We propose two complementary solutions.
On the SSM side, \textbf{SM-SSM} (Selectivity-Modulated SSM) smoothly interpolates between selective Mamba dynamics and fixed LTI dynamics (mathematically equivalent to a learned causal convolution) via an SNR-derived gate, with provable superset guarantee and $O(\sigma_n^6) \to O(\sigma_n^2)$ variance reduction (+38 parameters).
We extend this to \textbf{NC-SSM} (Noise-Conditioned SM-SSM), which replaces scalar SNR gating with per-sub-band selectivity---mapping each SSM state to a frequency sub-band---plus a \textbf{Learned Spectral Gate (LSG)} for explicit input denoising.
On the CNN side, \textbf{SAGN} (SNR-Adaptive Gated Network) augments BC-ResNet with the same LSG front-end and per-layer \textbf{SNR-Conditioned Channel Scale (SCCS)}, patching CNN's input-agnostic processing while preserving its structural noise advantages.
We further decompose BC-ResNet-1's noise robustness into four structural mechanisms (fixed kernels, Sub-Spectral Norm, broadcast residual, progressive abstraction), identifying three exploitable weaknesses.
Under capacity-controlled comparison ($\leq$7{,}464 parameters), all models are evaluated on Google Speech Commands V2 across 5~noise types $\times$ 7~SNR levels.
% TODO: Add results summary after training.
\end{abstract}

\begin{IEEEkeywords}
Keyword spotting, state space model, selectivity modulation, noise robustness, SNR-adaptive, edge AI
\end{IEEEkeywords}


% ============================================================
% I. INTRODUCTION
% ============================================================
\section{Introduction}
\label{sec:intro}

\IEEEPARstart{K}{eyword} spotting (KWS) enables always-on voice interfaces on microcontrollers (MCUs) and edge devices, where model size is constrained to sub-256\,KB and inference latency must remain below 10\,ms~\cite{banbury2021mlperf,lin2020mcunet}.
State space models (SSMs)~\cite{gu2024mamba,gu2022efficiently}, particularly the selective SSM architecture Mamba~\cite{gu2024mamba}, offer linear-time sequential modeling with input-dependent selection, making them attractive for edge deployment due to constant-memory streaming inference.
However, a persistent and puzzling gap remains: \emph{under noisy conditions with noise-augmented training, CNN architectures consistently outperform SSMs of comparable size}.

For example, BC-ResNet-1~\cite{kim2021bcresnet} (7{,}464 parameters) outperforms NanoMamba (7{,}402 parameters) under noise augmentation, despite nearly identical parameter budgets.
Conversely, when both are trained on clean data only and evaluated under noise, the SSM retains substantially higher accuracy than CNNs at moderate SNR ($-5$ to $15$\,dB), demonstrating that SSMs possess \emph{structural} noise robustness that CNNs lack---yet this advantage is \emph{reversed} after noise-augmented training.

In this work, we provide a unified analysis of this phenomenon from both the SSM and CNN perspectives, and propose targeted architectural solutions for each paradigm.
We identify \textbf{multiplicative noise propagation} as the structural cause of SSM degradation: input-dependent dynamics ($\Delta$, $\mathbf{B}$, $\mathbf{C}$ all derived from $x$) create a cubic noise coupling $\Delta(x) \cdot B(x) \cdot x$, yielding $O(\sigma_n^6)$ output variance versus $O(\sigma_n^2)$ for CNN fixed kernels.
Simultaneously, we decompose CNN noise robustness into four structural mechanisms---fixed kernels, Sub-Spectral Normalization, broadcast residuals, and progressive abstraction---while identifying three exploitable weaknesses: input-agnostic processing, no explicit denoising, and uniform frequency treatment.

Our contributions are:
\begin{enumerate}
\setlength\itemsep{0pt}
\item \textbf{Noise propagation analysis}: Formal characterization of the multiplicative ($O(\sigma_n^6)$) vs.\ additive ($O(\sigma_n^2)$) noise variance gap between selective SSMs and CNNs, with proof of LTI-SSM $\equiv$ learned causal convolution equivalence (Section~\ref{sec:analysis}).

\item \textbf{SM-SSM}: Selectivity-Modulated SSM that smoothly interpolates between selective (Mamba) and fixed (LTI) dynamics via an SNR-derived gate $\sigma$, with superset guarantee and $O(\sigma_n^6) \to O(\sigma_n^2)$ variance reduction, adding only +38~parameters (Section~\ref{sec:smssm}). Extended from~\cite{choi2026nanomamba_interspeech}.

\item \textbf{NC-SSM}: Noise-Conditioned SM-SSM extending SM-SSM with per-sub-band selectivity (5~sub-bands $\to$ 5~SSM states), stationarity-conditioned $\Delta$-floor, spectral-flatness-conditioned $\mathbf{B}$-base, and Learned Spectral Gate (+142~parameters; Section~\ref{sec:ncssm}).

\item \textbf{SAGN}: SNR-Adaptive Gated Network augmenting BC-ResNet with LSG front-end and per-layer SNR-Conditioned Channel Scale, patching CNN weaknesses while preserving structural advantages (Section~\ref{sec:sagn}).

\item \textbf{CNN structural analysis}: Decomposition of BC-ResNet-1's noise robustness into four mechanisms and three exploitable weaknesses (Section~\ref{sec:cnn_analysis}).

\item \textbf{Capacity-controlled evaluation}: Comprehensive comparison of 6~models (all $\leq$7{,}464 parameters) across 5~noise types $\times$ 7~SNR levels on GSC~V2 (Section~\ref{sec:results}).
\end{enumerate}


% ============================================================
% II. RELATED WORK
% ============================================================
\section{Related Work}
\label{sec:related}

\subsection{Small-Footprint Keyword Spotting}

The dominant paradigm for edge KWS uses CNNs on mel-spectrogram features.
DS-CNN~\cite{zhang2017hello} introduced depthwise separable convolutions, achieving 95.4\% on GSC with 23.7K parameters.
BC-ResNet~\cite{kim2021bcresnet} improved efficiency through broadcasted residual connections and Sub-Spectral Normalization (SSN), achieving 96.0\% with 7.5K parameters.
TC-ResNet~\cite{choi2019temporal} proposed temporal convolutions for real-time inference, while MatchboxNet~\cite{majumdar2020matchboxnet} combined 1D separable convolutions with squeeze-and-excitation.
Attention-based approaches~\cite{deandrade2018neural,berg2021keyword} achieve strong accuracy but with $O(T^2)$ complexity.
Recent system-level work~\cite{bartoli2025typman} has demonstrated end-to-end KWS deployment on embedded MCUs.
A common limitation is that noise robustness depends on training-time BatchNorm~\cite{ioffe2015batchnorm} statistics, which become mismatched under noise distribution shift at inference.

\subsection{State Space Models for Audio}

Structured State Spaces (S4)~\cite{gu2022efficiently} established continuous-time dynamical systems with HiPPO~\cite{gu2020hippo} initialization as effective for long-range modeling.
Mamba~\cite{gu2024mamba} introduced input-dependent \emph{selective} dynamics, and Mamba-2~\cite{dao2024mamba2} unified SSMs with attention.
For audio, Keyword Mamba~\cite{goel2024keyword} applied bidirectional Mamba to KWS but required 3.4M parameters.
SSAMBA~\cite{shaker2024ssamba} explored self-supervised audio representations.
Our prior SA-SSM~\cite{choi2026nanomamba_interspeech} was the first to condition SSM dynamics ($\Delta$, $\mathbf{B}$) on per-band SNR for noise-adaptive temporal processing, and SM-SSM introduced the selectivity modulation mechanism.
Multi-axis SSM approaches (Vision Mamba~\cite{zhu2024vision}, VMamba~\cite{liu2024vmamba}) scan 2D patches in multiple directions but do not address audio-specific noise robustness.

\subsection{Noise-Robust KWS}

Multi-condition training~\cite{ko2015audio} and SpecAugment~\cite{park2019specaugment} are standard training-time approaches for noise robustness.
Classical signal processing methods---spectral subtraction~\cite{boll1979suppression}, Wiener filtering~\cite{haykin2014adaptive}, and minimum mean-square error (MMSE) estimation~\cite{ephraim1984speech}---provide explicit noise removal but introduce artifacts.
PCEN~\cite{wang2017trainable} replaces static log-mel with trainable per-channel energy normalization for robust front-ends.
An overview of noise-robust ASR techniques is provided by~\cite{li2014overview}.
Our work is the first to analyze the SSM--CNN noise robustness gap from both theoretical (noise variance analysis) and architectural (structural mechanism decomposition) perspectives.

\subsection{SNR-Adaptive and Dynamic Networks}

Dynamic neural networks~\cite{han2021dynamic} adapt computation based on input characteristics.
Noise-conditioned mixture of experts~\cite{zhang2025noisemoe} route computation based on noise estimates.
Our SM-SSM, NC-SSM, and SAGN fall within this paradigm: all three condition their processing on SNR estimates, but through distinct mechanisms---selectivity modulation (SM-SSM/NC-SSM) versus channel scaling (SAGN).


% ============================================================
% III. NOISE PROPAGATION ANALYSIS
% ============================================================
\section{Noise Propagation Analysis}
\label{sec:analysis}

This section provides the theoretical foundation for all subsequent architectural proposals, formalizing the structural difference between SSM and CNN noise behavior.

\subsection{Selective SSM: Multiplicative Noise Coupling}
\label{sec:mult_noise}

A selective SSM~\cite{gu2024mamba} processes input $x_t \in \mathbb{R}^D$ via:
\begin{align}
[\Delta_t, \mathbf{B}_t, \mathbf{C}_t] &= W_{\text{proj}} \cdot x_t, \label{eq:proj} \\
\mathbf{h}_t &= \bar{\mathbf{A}}_t \mathbf{h}_{t-1} + \bar{\mathbf{B}}_t x_t, \quad
y_t = \mathbf{C}_t \mathbf{h}_t + \mathbf{D} x_t, \label{eq:ssm}
\end{align}
where $\bar{\mathbf{A}}_t = \exp(\mathbf{A} \cdot \Delta_t)$, $\bar{\mathbf{B}}_t = \Delta_t \cdot \mathbf{B}_t$.
The input contribution to the state update is:
\begin{equation}
u_t = \underbrace{\Delta_t}_{\text{from } x} \cdot \underbrace{\mathbf{B}_t}_{\text{from } x} \cdot \underbrace{x_t}_{\text{input}} = f_\Delta(x_t) \cdot f_B(x_t) \cdot x_t.
\label{eq:mult}
\end{equation}
All three factors depend on $x_t$. When $x_t = s_t + n_t$:
\begin{equation}
u_t = \underbrace{f_\Delta(s) f_B(s) s}_{\text{clean}} + \underbrace{f_\Delta(n) f_B(n) n}_{\text{noise}^3} + \text{mixed terms}.
\label{eq:cubic}
\end{equation}

\begin{proposition}[Noise variance---selective SSM]
\label{prop:var_ssm}
Let $n \sim \mathcal{N}(0, \sigma_n^2 I)$. For a selective SSM with linear projections $f_\Delta$, $f_B$:
\begin{equation}
\mathrm{Var}[u_t] = O(\sigma_n^6).
\label{eq:var_ssm}
\end{equation}
\end{proposition}
\begin{proof}
Each of the three factors in~\eqref{eq:mult} contributes $O(\sigma_n^2)$ variance through linear projections of~$n$. Since the factors are multiplied, the dominant noise-only term has variance $O(\sigma_n^{2 \times 3}) = O(\sigma_n^6)$.
\end{proof}


\subsection{CNN / LTI-SSM: Additive Noise}
\label{sec:add_noise}

A CNN convolves input with \emph{fixed, learned} kernels $W$:
\begin{equation}
y = W \ast (s + n) = \underbrace{W \ast s}_{\text{signal}} + \underbrace{W \ast n}_{\text{noise (additive)}}.
\label{eq:cnn}
\end{equation}
Since $W$ is independent of input, noise enters only through a single linear transformation.

\begin{proposition}[Noise variance---CNN/LTI]
\label{prop:var_cnn}
For a CNN or LTI-SSM with fixed parameters $\theta_\Delta$, $\bm{\theta}_B$:
\begin{equation}
\mathrm{Var}[u_t] = \theta_\Delta^2 \lVert\bm{\theta}_B\rVert^2 \sigma_n^2 = O(\sigma_n^2).
\label{eq:var_cnn}
\end{equation}
\end{proposition}

\smallskip
The variance ratio at $-15$\,dB ($\sigma_n \approx 5.6$) is $\sigma_n^6 / \sigma_n^2 = \sigma_n^4 \approx 1{,}000$: selective SSM output variance is \emph{three orders of magnitude} higher than LTI/CNN.


\subsection{LTI-SSM as Learned Causal Convolution}
\label{sec:lti_equiv}

An LTI SSM uses \emph{fixed} $\bar{A}$, $\bar{B}$, $C$. Unrolling the recurrence:
\begin{equation}
y_t = \sum_{k=0}^{t} \underbrace{C \bar{A}^{t-k} \bar{B}}_{K_{t-k}} x_k + D x_t = (K \ast x)_t + D x_t,
\label{eq:conv}
\end{equation}
where $K_\tau = C \bar{A}^\tau \bar{B}$ is a fixed, exponentially decaying causal IIR kernel:
\begin{equation}
\boxed{\text{LTI-SSM} \equiv \text{Causal IIR Convolution} \equiv \text{Learned Conv1D}.}
\label{eq:equivalence}
\end{equation}
This establishes that transitioning from selective to LTI dynamics at low SNR is mathematically equivalent to replacing input-dependent processing with a \emph{fixed learned convolution}---achieving CNN-like structural noise immunity.


\subsection{CNN Structural Advantage Decomposition}
\label{sec:cnn_analysis}

We decompose BC-ResNet-1's noise robustness into four structural mechanisms that explain why it outperforms SSMs after noise-augmented training:

\smallskip\noindent\textbf{1) Fixed kernels as noise-tolerant feature detectors.}
Under noise augmentation, each convolutional kernel optimizes for noisy speech patterns. Since kernels are input-independent, they apply the same stable feature extraction regardless of input SNR. Noise propagates additively ($O(\sigma_n^2)$), making the learned filters structurally stable across SNR conditions.

\smallskip\noindent\textbf{2) Sub-Spectral Normalization (SSN).}
BC-ResNet divides 40 mel bins into sub-bands (e.g., 5 groups of 8 bins) and applies separate BatchNorm per sub-band. This \emph{structurally} equalizes frequency-dependent noise energy: factory noise concentrated in low-frequency bins is normalized independently from high-frequency bins, preventing low-band noise from dominating statistics.

\smallskip\noindent\textbf{3) Broadcast residual connections.}
BC-ResNet pools across frequency ($\text{freq\_pool}: (B,C,F,T) \to (B,C,1,T)$) then broadcasts back. This averages noise energy across all frequency bands: if noise concentrates in specific bands, broadcasting \emph{dilutes} it across all bands, providing frequency-agnostic noise smoothing.

\smallskip\noindent\textbf{4) Deep progressive abstraction.}
Seven residual blocks progressively transform spectral features: early layers capture low-level spectral patterns (noise-sensitive), middle layers capture phoneme-like patterns (moderately robust), and deep layers capture keyword-level patterns (highly robust). Each layer abstracts away noise through its fixed learned filters.


\subsection{CNN Weaknesses: Design Implications}
\label{sec:cnn_weakness}

Despite these strengths, CNN architectures have three exploitable weaknesses:

\smallskip\noindent\textbf{W1: Input-agnostic processing.}
The same frozen kernels and BatchNorm statistics apply to both clean signals and $-15$\,dB signals identically. The CNN cannot know ``this input has low SNR'' and adjust accordingly.

\smallskip\noindent\textbf{W2: No explicit denoising.}
CNNs learn to \emph{ignore} noise through invariant features, but never explicitly \emph{remove} it. At low SNR, the noise energy exceeds what can be ignored.

\smallskip\noindent\textbf{W3: Uniform frequency treatment.}
While SSN normalizes sub-bands separately, the convolution kernels apply identically across all frequencies. Noisy bands are not selectively suppressed.

\smallskip
These weaknesses motivate SAGN (Section~\ref{sec:sagn}), which patches W1 with SNR-Conditioned Channel Scale and W2/W3 with Learned Spectral Gate. On the SSM side, the multiplicative noise propagation motivates SM-SSM (Section~\ref{sec:smssm}) and NC-SSM (Section~\ref{sec:ncssm}).


% ============================================================
% IV. SM-SSM: SELECTIVITY-MODULATED SSM
% ============================================================
\section{SM-SSM: Selectivity-Modulated SSM}
\label{sec:smssm}

SM-SSM introduces an SNR-adaptive gate $\sigma$ that blends selective (input-dependent) and fixed (LTI) dynamics, realizing a universal principle: \emph{selectivity should be inversely proportional to noise level}.

\subsection{Selectivity Modulation Formulation}

Three enhancements address limitations of naive selectivity modulation:

\smallskip\noindent\textbf{(i) Temporal $\sigma$ smoothing.}
Raw per-frame SNR estimates are noisy, causing gate instability. We apply a causal 3-frame average:
$\tilde{s}_t = \frac{1}{3}\sum_{k=0}^{2} \bar{s}_{t-k}$, where $\bar{s}_t = \text{mean}_f\big(s_f/(s_f{+}K)\big)$ uses Michaelis-Menten normalization ($K{=}0.05$).

\smallskip\noindent\textbf{(ii) Per-state selectivity.}
$\Delta$ and $(\mathbf{B}, \mathbf{C})$ use separate learnable biases:
\begin{align}
\sigma^\Delta_t &= \text{sigmoid}(s_{\text{scale}} \cdot \tilde{s}_t + b_\Delta), \label{eq:sigma_dt} \\
\sigma^{BC}_{t,n} &= \text{sigmoid}(s_{\text{scale}} \cdot \tilde{s}_t + b_{BC,n}), \label{eq:sigma_bc}
\end{align}
where $b_\Delta \in \mathbb{R}$ and $\mathbf{b}_{BC} \in \mathbb{R}^N$, enabling some states to maintain selectivity while others transition to LTI for noise immunity.

\smallskip\noindent\textbf{(iii) PCEN-conditioned $\sigma$.}
For non-stationary noise, $\sigma$ is modulated by the DualPCEN routing gate:
$\sigma \leftarrow \sigma \cdot (1 - \lambda_g \cdot g_t)$,
where $\lambda_g$ is learnable ($\lambda_g{=}0.3$ init).

\smallskip\noindent\textbf{Blending equations:}
\begin{align}
\Delta_t &= \sigma^\Delta_t \cdot \Delta_{\text{sel}}(x_t) + (1{-}\sigma^\Delta_t) \cdot \theta_\Delta, \label{eq:blend_dt} \\
\mathbf{B}_t &= \bm{\sigma}^{BC}_t \odot \mathbf{B}_{\text{sel}}(x_t) + (1{-}\bm{\sigma}^{BC}_t) \odot \bm{\theta}_B, \label{eq:blend_b} \\
\mathbf{C}_t &= \bm{\sigma}^{BC}_t \odot \mathbf{C}_{\text{sel}}(x_t) + (1{-}\bm{\sigma}^{BC}_t) \odot \bm{\theta}_C, \label{eq:blend_c}
\end{align}
where $\theta_\Delta$, $\bm{\theta}_B$, $\bm{\theta}_C$ are learned fixed (LTI) parameters.

\smallskip\noindent\textbf{Two extremes:}
$\sigma \approx 1$ (high SNR): full selective SSM (Mamba).
$\sigma \approx 0$ (low SNR): LTI-SSM $\equiv$ learned convolution~\eqref{eq:equivalence}.


\subsection{Theoretical Guarantees}

\begin{proposition}[Superset guarantee]
\label{prop:superset}
Let $\mathcal{H}_{\text{SSM}}$ denote the hypothesis space of the standard selective SSM, and $\mathcal{H}_{\text{SM}}$ that of SM-SSM. Then $\mathcal{H}_{\text{SM}} \supset \mathcal{H}_{\text{SSM}}$.
\end{proposition}
\begin{proof}
Setting $\sigma{=}1$ identically recovers $\Delta_t = \Delta_{\text{sel}}(x_t)$, $\mathbf{B}_t = \mathbf{B}_{\text{sel}}(x_t)$, $\mathbf{C}_t = \mathbf{C}_{\text{sel}}(x_t)$---the standard SSM.
Since SM-SSM additionally includes all $\sigma{<}1$ configurations, $\mathcal{H}_{\text{SM}} \supset \mathcal{H}_{\text{SSM}}$.
\end{proof}

This ensures $\mathcal{L}^*_{\text{SM}} \leq \mathcal{L}^*_{\text{SSM}}$: SM-SSM can only improve, never degrade.

\begin{proposition}[Variance reduction at low SNR]
\label{prop:var_reduction}
At low SNR where $\sigma \approx 0$, SM-SSM reduces state update variance from $O(\sigma_n^6)$~\eqref{eq:var_ssm} to $O(\sigma_n^2)$~\eqref{eq:var_cnn}.
\end{proposition}
\begin{proof}
At $\sigma{=}0$: $\Delta_t = \theta_\Delta$, $\mathbf{B}_t = \bm{\theta}_B$ (fixed). Thus $u_t = \theta_\Delta \cdot \bm{\theta}_B \cdot x_t$, where only $x_t$ depends on noise, matching the LTI/CNN case.
\end{proof}


\subsection{Parameter Overhead}

SM-SSM adds $3N{+}4$ parameters per block: $\theta_\Delta$ (1), $\bm{\theta}_B$ ($N$), $\bm{\theta}_C$ ($N$), $s_{\text{scale}}$ (1), $b_\Delta$ (1), $\mathbf{b}_{BC}$ ($N$), $\lambda_g$ (1).
With $N{=}5$ and 2 blocks: \textbf{+38 parameters total (0.51\%)}.
Initialization ($\theta_\Delta{=}0$, $\bm{\theta}_B{=}\mathbf{0}$, $s_{\text{bias}}{=}{-1}$ yielding $\sigma{\approx}0.97$ at clean) ensures training starts functionally identical to the standard SSM.


% ============================================================
% V. NC-SSM: NOISE-CONDITIONED SM-SSM
% ============================================================
\section{NC-SSM: Noise-Conditioned SM-SSM}
\label{sec:ncssm}

SM-SSM uses a scalar mean SNR to derive $\sigma$, losing frequency-dependent information. NC-SSM extends SM-SSM with three innovations.

\subsection{Limitation of Scalar Selectivity}

SM-SSM computes a single scalar $\bar{s}_t = \text{mean}_f(\hat{s}_{f,t})$ from 40 mel-band SNR estimates, then derives one $\sigma^\Delta$ and per-state $\sigma^{BC}_n$ values all from this same scalar.
Under factory noise (energy concentrated below 500\,Hz), the mean SNR is low even though high-frequency bands are clean.
Consequently, \emph{all} SSM states transition toward LTI, unnecessarily sacrificing selectivity in clean frequency regions.

\subsection{Per-Sub-Band Selectivity}
\label{sec:subband}

NC-SSM pools 40 mel bands into $N{=}5$ sub-bands of 8 bands each, and maps each sub-band to one SSM state:
\begin{equation}
\bar{s}^{(n)}_t = \frac{1}{8} \sum_{f=8n}^{8n+7} \hat{s}_{f,t}, \quad n = 0, \ldots, 4.
\label{eq:subband_snr}
\end{equation}

The per-state selectivity gate becomes:
\begin{equation}
\sigma^{BC}_{t,n} = \text{sigmoid}\big(w^{\text{sub}}_n \cdot \bar{s}^{(n)}_t + b_{BC,n}\big),
\label{eq:nc_sigma}
\end{equation}
where $\mathbf{w}^{\text{sub}} \in \mathbb{R}^N$ replaces the single $s_{\text{scale}}$ for $\mathbf{B}$/$\mathbf{C}$ modulation.
The temporal gate $\sigma^\Delta$ remains scalar (controls global temporal memory).

\smallskip
Under factory noise:
\begin{itemize}
\setlength\itemsep{0pt}
\item State 0 (low-freq, 0--500\,Hz): $\sigma \approx 0$ (LTI, noise defense)
\item States 3--4 (high-freq, 3--8\,kHz): $\sigma \approx 1$ (selective, clean signal)
\end{itemize}
This preserves selective dynamics where the signal is clean while protecting noisy bands---impossible with SM-SSM's scalar approach.


\subsection{Stationarity-Conditioned $\Delta$ Floor}
\label{sec:station}

The DualPCEN routing gate $g_t$ indicates noise stationarity (high $g_t$ = stationary). NC-SSM modulates the $\Delta$-floor:
\begin{equation}
\Delta_{\text{floor},t} = \Delta_{\text{base}} + \alpha_s \cdot g_t,
\label{eq:station_dt}
\end{equation}
where $\alpha_s$ is learnable (init.~0).
Under stationary noise, a higher $\Delta$-floor enforces longer temporal memory for noise averaging.
Under non-stationary noise (babble), a lower floor allows faster adaptation.

\subsection{Spectral-Flatness-Conditioned $\mathbf{B}$ Base}
\label{sec:sf_base}

NC-SSM detects broadband vs.\ narrowband noise via a broadband score:
\begin{equation}
\beta_t = 1 - \frac{\text{std}(\bar{\mathbf{s}}_t)}{\text{mean}(\bar{\mathbf{s}}_t) + \epsilon},
\label{eq:broadband}
\end{equation}
where $\bar{\mathbf{s}}_t$ is the 5-element sub-band SNR vector.
Broadband noise (white, pink) yields $\beta_t \approx 1$ (uniform SNR); narrowband (factory) yields $\beta_t \approx 0$.
The LTI-mode input coupling is modulated:
\begin{equation}
\tilde{\bm{\theta}}_B = \bm{\theta}_B \cdot (1 + \mathbf{w}_{sf} \odot \beta_t),
\label{eq:sf_B}
\end{equation}
where $\mathbf{w}_{sf} \in \mathbb{R}^N$ (init.~$\mathbf{0}$) strengthens fixed input coupling under broadband noise.


\subsection{Learned Spectral Gate (LSG)}
\label{sec:lsg}

Both NC-SSM and SAGN share a \textbf{Learned Spectral Gate}---a lightweight per-frequency, per-frame gating module applied before feature normalization:
\begin{equation}
\text{LSG}(\mathbf{m}_t, \hat{\mathbf{s}}_t) = \mathbf{m}_t \odot \big(\mathbf{g}_t + (1{-}\mathbf{g}_t) \odot \bm{\phi}\big),
\label{eq:lsg}
\end{equation}
where:
\begin{equation}
\mathbf{g}_t = \text{sigmoid}(\mathbf{w} \odot \hat{\mathbf{s}}_t + \mathbf{b}), \quad
\bm{\phi} = \text{sigmoid}(\mathbf{f}),
\label{eq:lsg_gate}
\end{equation}
with learnable $\mathbf{w}, \mathbf{b}, \mathbf{f} \in \mathbb{R}^F$ (total: $3F = 120$ parameters for $F{=}40$).
At high SNR, $\mathbf{g} \approx 1$ (pass-through); at low SNR, $\mathbf{g} \approx 0$ and $\bm{\phi}$ provides a learned frequency-dependent floor, suppressing noisy bands while preserving a minimum signal level.

LSG implements \emph{learned spectral gating}~\cite{ephraim1984speech} end-to-end, providing explicit per-band denoising that neither vanilla SSM nor CNN architectures possess.

\subsection{Warm-Start Initialization}

All NC-SSM-specific parameters are initialized for exact SM-SSM equivalence:
$\mathbf{w}^{\text{sub}} = 5.0$ (matching $s_{\text{scale}}$),
$\alpha_s = 0$ (no stationarity effect),
$\mathbf{w}_{sf} = \mathbf{0}$ (no spectral-flatness effect).
This enables NC-SSM to start from SM-SSM's learned behavior and fine-tune the new parameters.


\subsection{NC-SSM Parameter Summary}

NC-SSM adds per block: $\mathbf{w}^{\text{sub}}$ ($N$), $\alpha_s$ (1), $\mathbf{w}_{sf}$ ($N$) = $2N{+}1 = 11$ parameters.
With 2 blocks: +22.
LSG: +120.
\textbf{Total NC-SSM model: 7{,}171 parameters} (269 fewer than BC-ResNet-1).


% ============================================================
% VI. SAGN: SNR-ADAPTIVE GATED NETWORK
% ============================================================
\section{SAGN: SNR-Adaptive Gated Network}
\label{sec:sagn}

Rather than only fixing SSM weaknesses, SAGN fixes CNN weaknesses. It takes BC-ResNet-1's proven backbone and adds two modules targeting the weaknesses identified in Section~\ref{sec:cnn_weakness}.

\subsection{Architecture}

SAGN's processing pipeline is:
\begin{equation}
\text{Audio} \xrightarrow{\text{STFT}} \text{Mel} \xrightarrow{\text{LSG}} \hat{\mathbf{M}} \xrightarrow{\text{IN}} \tilde{\mathbf{M}} \xrightarrow{\text{BC-ResNet}} \xrightarrow{\text{GAP+FC}} \hat{y},
\label{eq:sagn_pipeline}
\end{equation}
where LSG provides explicit denoising (addresses W2, W3 from Section~\ref{sec:cnn_weakness}), and each BC-ResNet stage is followed by SNR-Conditioned Channel Scale (addresses W1).

The BC-ResNet backbone uses identical structure to BC-ResNet-1: initial Conv2d$(1, 8, 5{\times}5)$, then 4 stages with channel progression $8 \to 8 \to 12 \to 16 \to 20$ using BCResBlocks with Sub-Spectral Normalization and broadcast residuals.

\subsection{LSG Front-End}

The same Learned Spectral Gate module from Section~\ref{sec:lsg} (120 parameters) is applied to the mel spectrogram before InstanceNorm.
In the CNN context, LSG replaces the lack of explicit denoising: rather than relying solely on learned filter invariance, SAGN first reduces noise energy at the input level.

\subsection{SNR-Conditioned Channel Scale (SCCS)}
\label{sec:sccs}

After each BC-ResNet stage, SCCS applies per-channel scaling conditioned on global SNR:
\begin{equation}
\text{SCCS}(\mathbf{x}, \bar{s}) = \mathbf{x} \odot (\bm{\beta} + \bm{\alpha} \cdot \bar{s}),
\label{eq:sccs}
\end{equation}
where $\bar{s} = \text{mean}(\hat{\mathbf{s}})$ is the global SNR estimate, $\bm{\alpha}, \bm{\beta} \in \mathbb{R}^C$ are learnable (init.\ $\bm{\alpha}{=}\mathbf{0}$, $\bm{\beta}{=}\mathbf{1}$), and $C$ is the number of channels in that stage.

At initialization, SCCS is an identity operation. During training, it learns to:
\begin{itemize}
\setlength\itemsep{0pt}
\item At low SNR ($\bar{s} \approx 0$): suppress noise-sensitive channels via reduced $\bm{\beta}$.
\item At high SNR ($\bar{s} \approx 1$): fully utilize all channels via $\bm{\beta} + \bm{\alpha}$.
\end{itemize}

SCCS adds $2C$ parameters per stage: $2 \times (8 + 12 + 16 + 20) = 112$ parameters across 4 stages.


\subsection{Design Philosophy}

SAGN preserves all four CNN structural advantages (Section~\ref{sec:cnn_analysis}) while patching the three weaknesses:
\begin{itemize}
\setlength\itemsep{0pt}
\item W1 (input-agnostic) $\to$ SCCS makes each layer SNR-aware
\item W2 (no denoising) $\to$ LSG provides explicit spectral gating
\item W3 (uniform frequency) $\to$ LSG gates per-frequency band based on SNR
\end{itemize}

The same SNR estimator and LSG module are shared between NC-SSM and SAGN, enabling fair comparison: both paradigms benefit from the same front-end innovation.

\textbf{Total SAGN: 7{,}240 parameters} (224 fewer than BC-ResNet-1).


% ============================================================
% VII. EXPERIMENTAL SETUP
% ============================================================
\section{Experimental Setup}
\label{sec:setup}

\subsection{Dataset}

We evaluate on Google Speech Commands V2~\cite{warden2018speech} with the standard 12-class task (10 keywords + silence + unknown): 86{,}843 training / 10{,}481 validation / 11{,}505 test utterances at 16\,kHz.

\subsection{Noise Conditions}

Five noise types are evaluated:
\textbf{factory} (machine hum at 50--250\,Hz harmonics, conveyor rumble, impact transients),
\textbf{white} (Gaussian, spectrally flat),
\textbf{babble} (5--9 randomly mixed training utterances),
\textbf{street} (traffic rumble, horn impulses, road noise), and
\textbf{pink} ($1/f$ spectrum).
Noise is mixed at target SNR using RMS-based scaling at levels $\{-15, -10, -5, 0, 5, 10, 15\}$\,dB, yielding $5 \times 7 = 35$ noise conditions plus clean.

\subsection{Models Compared}

\begin{table}[t]
\centering
\caption{Model configurations. All primary models (except DS-CNN-S) are within BC-ResNet-1's 7{,}464-parameter budget for capacity-controlled comparison.}
\label{tab:configs}
\small
\setlength{\tabcolsep}{3pt}
\begin{tabular}{llrc}
\toprule
\textbf{Model} & \textbf{Type} & \textbf{Params} & \textbf{INT8} \\
\midrule
DS-CNN-S~\cite{zhang2017hello} & CNN & 23{,}756 & 23.2\,KB \\
BC-ResNet-1~\cite{kim2021bcresnet} & CNN & 7{,}464 & 7.3\,KB \\
\midrule
NM-Matched (SA-SSM v2) & SSM & 7{,}402 & 7.2\,KB \\
SM-SSM & SSM & 7{,}440 & 7.3\,KB \\
\textbf{NC-SSM} & SSM+LSG & \textbf{7{,}171} & \textbf{7.0\,KB} \\
\textbf{SAGN} & CNN+LSG+SCCS & \textbf{7{,}240} & \textbf{7.1\,KB} \\
\bottomrule
\end{tabular}
\end{table}

Table~\ref{tab:configs} lists all models. The primary comparison is between six models, all within BC-ResNet-1's 7{,}464-parameter budget (except DS-CNN-S as a larger reference). NC-SSM and SAGN are the most parameter-efficient at 7{,}171 and 7{,}240 respectively.

\subsection{Training Details}

All models are trained for 30 epochs with AdamW~\cite{loshchilov2019adamw} ($\beta_1{=}0.9$, $\beta_2{=}0.999$), cosine annealing~\cite{loshchilov2017sgdr} (initial LR $3{\times}10^{-3}$), label smoothing 0.1~\cite{szegedy2016rethinking}, batch size 128, and gradient clipping at 1.0.
Data augmentation includes time shift ($\pm$100\,ms), volume perturbation ($\pm$20\%), and light Gaussian noise ($p{=}0.3$, $\sigma{=}0.005$).

\textbf{Per-sample noise augmentation} with four-phase curriculum is applied to all models:
(1)~warm-up (epochs~0--2): clean only;
(2)~gentle (epochs~3--9): noise ratio ramps 0$\to$50\%, SNR~5--15\,dB;
(3)~moderate (epochs~10--19): full noise ratio, SNR~0--10\,dB;
(4)~hard (epochs~20+): full ratio, SNR~$-5$--10\,dB.
Each sample independently draws a noise type and SNR.
SSM models additionally use SpecAugment~\cite{park2019specaugment} and EMA (decay 0.999).


% ============================================================
% VIII. RESULTS
% ============================================================
\section{Results}
\label{sec:results}

All models are trained with per-sample noise augmentation and evaluated on the GSC~V2 12-class test set under 36 conditions (clean + 5 noise types $\times$ 7 SNR levels).

\subsection{Clean Accuracy}

\begin{table}[t]
\centering
\caption{Clean accuracy on GSC V2 (12-class). All models trained with per-sample noise augmentation and curriculum. $^\dagger$NC-SSM and SAGN are new contributions.}
\label{tab:clean}
\small
\begin{tabular}{lrr}
\toprule
\textbf{Model} & \textbf{Params} & \textbf{Val Acc (\%)} \\
\midrule
DS-CNN-S~\cite{zhang2017hello} & 23{,}756 & 96.4 \\
BC-ResNet-1~\cite{kim2021bcresnet} & 7{,}464 & 95.3 \\
\midrule
NM-Matched (SA-SSM v2) & 7{,}402 & 95.1 \\
SM-SSM & 7{,}440 & TBD \\
NC-SSM$^\dagger$ & 7{,}171 & TBD \\
SAGN$^\dagger$ & 7{,}240 & TBD \\
\bottomrule
\end{tabular}
\end{table}


\subsection{Noise Robustness}

Table~\ref{tab:noise_main} reports accuracy under five noise types at all SNR levels. Existing baseline results (NM-Matched, BC-ResNet-1, DS-CNN-S) are from prior training runs; SM-SSM, NC-SSM, and SAGN results are pending.

\begin{table*}[t]
\centering
\caption{Noise robustness comparison (accuracy \%). All models trained with per-sample noise augmentation. NM-Matched shows prior NanoMamba baseline. Best ultra-compact ($<$10K params) result per condition in \textbf{bold}. SM-SSM, NC-SSM, SAGN rows TBD pending training.}
\label{tab:noise_main}
\small
\begin{tabular}{ll|cccccccc}
\toprule
\textbf{Noise} & \textbf{Model} & \textbf{Clean} & $\mathbf{-15}$ & $\mathbf{-10}$ & $\mathbf{-5}$ & $\mathbf{0}$ & $\mathbf{5}$ & $\mathbf{10}$ & $\mathbf{15}$\,\textbf{dB} \\
\midrule
\multirow{6}{*}{Factory}
  & DS-CNN-S (23.8K)    & 96.8 & 64.4 & 74.6 & 85.9 & 91.5 & 93.8 & 95.1 & 95.5 \\
  & BC-ResNet-1 (7.5K)  & \textbf{95.3} & \textbf{64.5} & \textbf{74.8} & \textbf{83.5} & \textbf{89.0} & \textbf{91.9} & 92.9 & \textbf{93.9} \\
  & NM-Matched (7.4K)   & 95.1 & 61.1 & 70.6 & 82.3 & 88.5 & 91.4 & \textbf{93.0} & 93.8 \\
  & SM-SSM (7.4K)       & --- & --- & --- & --- & --- & --- & --- & --- \\
  & NC-SSM (7.2K)       & --- & --- & --- & --- & --- & --- & --- & --- \\
  & SAGN (7.2K)         & --- & --- & --- & --- & --- & --- & --- & --- \\
\midrule
\multirow{6}{*}{White}
  & DS-CNN-S (23.8K)    & 96.8 & 61.7 & 67.1 & 81.6 & 89.9 & 93.1 & 94.8 & 95.6 \\
  & BC-ResNet-1 (7.5K)  & \textbf{95.3} & \textbf{58.6} & \textbf{64.6} & \textbf{76.7} & \textbf{85.9} & \textbf{90.2} & 92.6 & 93.8 \\
  & NM-Matched (7.4K)   & 95.1 & 31.1 & 59.1 & 75.0 & 85.5 & 89.8 & \textbf{92.8} & \textbf{93.9} \\
  & SM-SSM (7.4K)       & --- & --- & --- & --- & --- & --- & --- & --- \\
  & NC-SSM (7.2K)       & --- & --- & --- & --- & --- & --- & --- & --- \\
  & SAGN (7.2K)         & --- & --- & --- & --- & --- & --- & --- & --- \\
\midrule
\multirow{6}{*}{Babble}
  & DS-CNN-S (23.8K)    & 96.8 & 79.2 & 86.2 & 91.4 & 94.5 & 95.5 & 95.9 & 96.5 \\
  & BC-ResNet-1 (7.5K)  & \textbf{95.3} & \textbf{76.7} & \textbf{84.3} & \textbf{88.8} & \textbf{92.1} & \textbf{93.8} & \textbf{94.4} & \textbf{95.0} \\
  & NM-Matched (7.4K)   & 95.1 & 71.4 & 79.0 & 86.0 & 91.0 & 93.0 & 93.9 & 94.5 \\
  & SM-SSM (7.4K)       & --- & --- & --- & --- & --- & --- & --- & --- \\
  & NC-SSM (7.2K)       & --- & --- & --- & --- & --- & --- & --- & --- \\
  & SAGN (7.2K)         & --- & --- & --- & --- & --- & --- & --- & --- \\
\midrule
\multirow{6}{*}{Street}
  & DS-CNN-S (23.8K)    & 96.8 & 61.1 & 69.4 & 81.5 & 88.9 & 92.8 & 95.0 & 95.9 \\
  & BC-ResNet-1 (7.5K)  & \textbf{95.3} & \textbf{60.7} & \textbf{65.1} & \textbf{77.7} & \textbf{84.8} & \textbf{90.1} & \textbf{93.0} & 94.1 \\
  & NM-Matched (7.4K)   & 95.1 & 57.6 & 64.2 & 73.9 & 83.8 & 88.4 & 92.4 & \textbf{94.2} \\
  & SM-SSM (7.4K)       & --- & --- & --- & --- & --- & --- & --- & --- \\
  & NC-SSM (7.2K)       & --- & --- & --- & --- & --- & --- & --- & --- \\
  & SAGN (7.2K)         & --- & --- & --- & --- & --- & --- & --- & --- \\
\midrule
\multirow{6}{*}{Pink}
  & DS-CNN-S (23.8K)    & 96.8 & 61.2 & 66.4 & 82.7 & 90.6 & 93.8 & 95.1 & 95.9 \\
  & BC-ResNet-1 (7.5K)  & \textbf{95.3} & \textbf{58.9} & \textbf{66.0} & \textbf{79.7} & \textbf{88.4} & \textbf{91.8} & \textbf{93.8} & \textbf{94.5} \\
  & NM-Matched (7.4K)   & 95.2 & 23.6 & 52.4 & 76.9 & 87.3 & 91.7 & 93.0 & 94.4 \\
  & SM-SSM (7.4K)       & --- & --- & --- & --- & --- & --- & --- & --- \\
  & NC-SSM (7.2K)       & --- & --- & --- & --- & --- & --- & --- & --- \\
  & SAGN (7.2K)         & --- & --- & --- & --- & --- & --- & --- & --- \\
\bottomrule
\end{tabular}
\end{table*}


\subsection{Structural vs.\ Training-Based Robustness}
\label{sec:struct_vs_train}

To disentangle \emph{structural} noise robustness (from architecture) from \emph{training-based} robustness (from noise data), we compare models trained on clean data only versus noise-augmented training, both evaluated at 0\,dB (Table~\ref{tab:clean_vs_aug}).

\begin{table}[t]
\centering
\caption{Clean-only vs.\ noise-aug training at 0\,dB. ``$\Delta$'' shows noise-aug improvement. SSMs demonstrate high structural robustness (small $\Delta$), confirmed at $-5$ to 15\,dB SNR range.}
\label{tab:clean_vs_aug}
\small
\begin{tabular}{llcccc}
\toprule
\textbf{Model} & \textbf{Train} & \textbf{Fact.} & \textbf{White} & \textbf{Bab.} & \textbf{Avg} \\
\midrule
\multirow{3}{*}{\shortstack[l]{NM-Matched\\(7.4K, SSM)}}
  & Clean  & 78.3 & 80.5 & 74.2 & 77.7 \\
  & NA     & 88.5 & 85.5 & 91.0 & 88.3 \\
  & $\Delta$ & \multicolumn{4}{c}{\textbf{+10.6\,\%p}} \\
\midrule
\multirow{3}{*}{\shortstack[l]{DS-CNN-S\\(23.8K, CNN)}}
  & Clean  & 75.6 & 13.9 & 70.1 & 53.2 \\
  & NA     & 91.5 & 89.9 & 94.5 & 92.0 \\
  & $\Delta$ & \multicolumn{4}{c}{+38.8\,\%p} \\
\midrule
\multirow{3}{*}{\shortstack[l]{BC-ResNet-1\\(7.5K, CNN)}}
  & Clean  & 71.6 & 54.7 & 73.7 & 66.7 \\
  & NA     & 89.0 & 85.9 & 92.1 & 89.0 \\
  & $\Delta$ & \multicolumn{4}{c}{+22.3\,\%p} \\
\midrule
\multirow{2}{*}{SM-SSM (7.4K)} & Clean / NA & \multicolumn{4}{c}{TBD} \\
\multirow{2}{*}{NC-SSM (7.2K)} & Clean / NA & \multicolumn{4}{c}{TBD} \\
\multirow{2}{*}{SAGN (7.2K)}   & Clean / NA & \multicolumn{4}{c}{TBD} \\
\bottomrule
\end{tabular}
\end{table}

The results reveal a striking asymmetry.
SSMs trained on clean data only retain \textbf{77.7\%} average accuracy at 0\,dB---11.0\,\%p higher than BC-ResNet-1 (\textbf{66.7\%}) and 24.5\,\%p higher than DS-CNN-S (\textbf{53.2\%}).
DS-CNN-S catastrophically collapses under white noise to 13.9\% (near random).
This confirms that SSM adaptive dynamics ($\Delta$-modulation, $\mathbf{B}$-gating) provide substantial \emph{inherent} noise robustness verified across $-5$ to 15\,dB SNR---without any noise during training.

However, after noise-augmented training, the ranking reverses: BC-ResNet-1 (89.0\%) surpasses NM-Matched (88.3\%).
This paradox---SSMs are structurally more robust but benefit less from noise augmentation---arises because CNN fixed kernels can fully absorb noise statistics during training ($O(\sigma_n^2)$ propagation), while SSM input-dependent parameters remain susceptible to multiplicative noise coupling ($O(\sigma_n^6)$) even after augmentation.

The noise-aug dependence ratio quantifies this: NM-Matched requires +10.6\,\%p from noise-aug, while DS-CNN-S requires +38.8\,\%p (3.7$\times$ more dependent on training data for noise robustness).


\subsection{Accuracy Drop Analysis}

\begin{table}[t]
\centering
\caption{Accuracy drop from clean to 0\,dB (\%p). Lower is better. SM-SSM, NC-SSM, SAGN: TBD.}
\label{tab:drop}
\small
\begin{tabular}{lccccc}
\toprule
\textbf{Model} & \textbf{Fact.} & \textbf{White} & \textbf{Bab.} & \textbf{Str.} & \textbf{Pink} \\
\midrule
DS-CNN-S    & \textbf{5.3} & \textbf{6.8} & \textbf{2.6} & \textbf{7.2} & \textbf{6.2} \\
BC-ResNet-1 & 5.9 & 9.6 & 3.4 & 10.0 & 7.1 \\
NM-Matched  & 6.6 & 9.6 & 4.1 & 11.3 & 7.9 \\
SM-SSM      & --- & --- & --- & --- & --- \\
NC-SSM      & --- & --- & --- & --- & --- \\
SAGN        & --- & --- & --- & --- & --- \\
\bottomrule
\end{tabular}
\end{table}


\subsection{Ablation Study}
\label{sec:ablation}

\begin{table}[t]
\centering
\caption{Ablation study. Evaluated at factory noise, 0\,dB. Each row adds one component. TBD after training.}
\label{tab:ablation}
\small
\begin{tabular}{lrcc}
\toprule
\textbf{Configuration} & \textbf{$\Delta$P} & \textbf{Clean} & \textbf{0\,dB} \\
\midrule
\multicolumn{4}{l}{\textit{SSM-side ablation:}} \\
Base SA-SSM v2         & 0     & 95.1 & 88.5 \\
\quad + SM-SSM modulation  & +38   & TBD  & TBD  \\
\quad + NC-SSM sub-band    & +22   & TBD  & TBD  \\
\quad + LSG front-end      & +120  & TBD  & TBD  \\
\quad + Station/SF cond.   & +0    & TBD  & TBD  \\
\midrule
\multicolumn{4}{l}{\textit{CNN-side ablation:}} \\
BC-ResNet-1 baseline   & 0     & 95.3 & 89.0 \\
\quad + LSG only           & +120  & TBD  & TBD  \\
\quad + SCCS only          & +112  & TBD  & TBD  \\
\quad + LSG + SCCS (SAGN)  & +232  & TBD  & TBD  \\
\bottomrule
\end{tabular}
\end{table}


\subsection{SS+Bypass Enhancement}
\label{sec:ss_bypass}

Spectral subtraction with an SNR-adaptive bypass mechanism provides additional improvement at extreme noise. Table~\ref{tab:ss_bypass} shows results at $-15$\,dB.

\begin{table}[t]
\centering
\caption{SS+Bypass improvement at $-15$\,dB SNR (\%p). Zero additional parameters. SM-SSM, NC-SSM, SAGN: TBD.}
\label{tab:ss_bypass}
\small
\begin{tabular}{lccccc}
\toprule
\textbf{Model} & \textbf{Fact.} & \textbf{White} & \textbf{Bab.} & \textbf{Str.} & \textbf{Pink} \\
\midrule
NM-Matched  & +3.6 & \textbf{+23.8} & $-$0.1 & +1.7 & \textbf{+29.2} \\
DS-CNN-S    & +1.3 & +1.3 & $-$0.1 & +2.9 & +1.8 \\
BC-ResNet-1 & +2.4 & +0.5 & $-$0.1 & $-$2.0 & +3.7 \\
SM-SSM      & --- & --- & --- & --- & --- \\
NC-SSM      & --- & --- & --- & --- & --- \\
SAGN        & --- & --- & --- & --- & --- \\
\bottomrule
\end{tabular}
\end{table}

Key findings from prior baselines:
(1)~SSMs benefit dramatically from SS+Bypass under stationary noise (+23.8\,\%p white, +29.2\,\%p pink at $-15$\,dB), while CNN baselines improve by only 0.5--3.7\,\%p.
(2)~All models show negligible change for babble ($-0.1$\,\%p), confirming the bypass gate correctly preserves original signals when SS would introduce artifacts.
(3)~Clean accuracy is preserved ($\pm$0\,\%p), validating the SNR-adaptive bypass design.


% ============================================================
% IX. ANALYSIS AND DISCUSSION
% ============================================================
\section{Analysis and Discussion}
\label{sec:discussion}

\subsection{Structural vs.\ Training-Based Robustness}

The clean-only vs.\ noise-aug comparison (Section~\ref{sec:struct_vs_train}) provides the clearest evidence for the distinction.
SSMs retain 77.7\% at 0\,dB with clean-only training, while DS-CNN-S collapses to 53.2\% (white: 13.9\%)---a 24.5\,\%p structural advantage despite 4.8$\times$ fewer parameters.
The noise-aug dependence ratio (SSM +10.6\,\%p vs DS-CNN-S +38.8\,\%p) directly quantifies this: CNN architectures are \emph{almost entirely dependent} on training data for noise robustness, while SSMs derive the majority from their architecture.

After noise-augmented training, CNN baselines achieve higher absolute accuracy due to their $O(\sigma_n^2)$ noise propagation---fixed kernels absorb noise statistics more efficiently during training.
SM-SSM and NC-SSM aim to reduce the SSM noise propagation from $O(\sigma_n^6)$ toward $O(\sigma_n^2)$, while SAGN aims to give CNNs the input-adaptive processing that SSMs naturally possess.

% TODO: Update with SM-SSM, NC-SSM, SAGN structural robustness results after training.

\subsection{Selectivity Gate Behavior}

% TODO: After SM-SSM and NC-SSM training:
% 1. Plot sigma vs SNR for SM-SSM (expected: sigmoid curve from ~1 at clean to ~0 at -15dB)
% 2. Plot per-sub-band sigma for NC-SSM under factory noise (expected: low-freq sigma near 0, high-freq near 1)
% 3. Compare SM-SSM (uniform sigma) vs NC-SSM (frequency-selective sigma)

\subsection{LSG as Architecture-Agnostic Innovation}

LSG provides explicit per-band, per-frame denoising using only 120 parameters, and is shared identically between NC-SSM and SAGN. This demonstrates that the innovation is \emph{architecture-agnostic}: both SSM and CNN pipelines benefit from learned spectral gating.

% TODO: Compare LSG contribution in NC-SSM context vs SAGN context from ablation results.

\subsection{When Does Each Architecture Win?}

% TODO: After results are available, analyze per-noise-type, per-SNR advantages:
% - High SNR (>10 dB): All models comparable; NC-SSM selectivity advantage?
% - Moderate SNR (0-10 dB): SAGN vs NC-SSM?
% - Low SNR (<0 dB): Type-dependent (factory: NC-SSM sub-band? white/pink: SAGN fixed kernels?)
% - Babble (non-stationary): PCEN-conditioned sigma helps NC-SSM?

\subsection{Deployment Considerations}

All models fit within MCU SRAM budgets of 64--256\,KB.
SSM models require fewer MACs due to linear-time processing: NanoMamba blocks use $O(L \cdot d_{\text{inner}} \cdot N)$ per time step vs.\ CNN's $O(K^2 \cdot C_{\text{in}} \cdot C_{\text{out}} \cdot F \cdot T)$.
SAGN, being CNN-based, benefits from mature deployment toolchains (TFLite Micro, CMSIS-NN).
INT8 quantization~\cite{rusci2020memory} is directly applicable to all models.
Runtime calibration (adjusting $\Delta$-floor, $\varepsilon$-bypass based on estimated noise level during silence) provides deployment-time adaptation for SSM models at zero parameter cost.

\subsection{Limitations and Future Work}

The current evaluation uses a single dataset (GSC V2) with procedurally generated noise.
Several directions remain:
(1)~real-world noise recordings;
(2)~larger-scale models (scaling to $\sim$24K params for DS-CNN-S comparison);
(3)~on-device latency and energy profiling on ARM Cortex-M4/M7;
(4)~hybrid architectures combining NC-SSM temporal processing with SAGN-like CNN frequency processing;
(5)~knowledge distillation~\cite{hinton2015distilling} from larger models.


% ============================================================
% X. CONCLUSION
% ============================================================
\section{Conclusion}
\label{sec:conclusion}

We identified \emph{multiplicative noise propagation} as the root cause of the selective SSM--CNN noise robustness gap: input-dependent dynamics create $O(\sigma_n^6)$ output variance versus $O(\sigma_n^2)$ for CNN fixed filters---a 1{,}000$\times$ gap at $-15$\,dB.
Rather than pursuing a single architectural paradigm, we proposed solutions from both sides.

On the SSM side, \textbf{SM-SSM} interpolates between selective Mamba and fixed LTI dynamics ($\equiv$ learned causal convolution) via an SNR-derived gate, with provable superset guarantee and variance reduction to $O(\sigma_n^2)$, adding only +38~parameters.
\textbf{NC-SSM} extends this with per-sub-band selectivity (each SSM state linked to a frequency sub-band) and Learned Spectral Gate for explicit denoising, totaling 7{,}171 parameters.

On the CNN side, \textbf{SAGN} preserves BC-ResNet-1's four structural advantages while patching three weaknesses through the same LSG front-end and SNR-Conditioned Channel Scale, totaling 7{,}240 parameters.

The shared LSG module (120 parameters) demonstrates that \emph{learned spectral gating is architecture-agnostic}: both SSM and CNN pipelines benefit from explicit per-band, per-frame SNR-conditioned denoising.

% TODO: Add experimental summary after results.

All models achieve these improvements within BC-ResNet-1's 7{,}464-parameter budget, demonstrating that the noise robustness gap is a \emph{structural} problem requiring \emph{structural} solutions---achievable with negligible parameter overhead.


\section*{Acknowledgments}
The code implementation, experimental analysis, and paper writing were assisted by Claude (Anthropic).
All research ideas, experimental design, and scientific decisions were made by the authors.


% ============================================================
% REFERENCES
% ============================================================
\bibliographystyle{IEEEtran}
\bibliography{refs_taslp}

\end{document}
